% Standalone sequel to Effective Domains, Oracles, and Forcing in the Effective Topos.
% Compile twice with pdfLaTeX. No external .bib, .aux, or figure files are required.
% References to Paper I use named results; its original source labels are recorded
% in Appendix A so that the sequel does not depend on an unstable theorem numbering.
\documentclass[11pt,a4paper]{article}
\usepackage[T1]{fontenc}
\usepackage[utf8]{inputenc}
\usepackage{lmodern}
\usepackage[margin=29mm]{geometry}
\usepackage{amsmath,amssymb,amsthm,mathtools}
\usepackage{microtype}
\usepackage{enumitem,booktabs,array,longtable}
\usepackage{tikz-cd}
\usepackage[colorlinks=true,linkcolor=blue!45!black,citecolor=blue!45!black,urlcolor=blue!45!black]{hyperref}
\usepackage{aliascnt}
\usepackage[nameinlink,noabbrev,capitalise]{cleveref}
\hypersetup{pdftitle={Effective Observation Hierarchies},pdfauthor={Conor Sheridan}}
\setlist{itemsep=2pt,topsep=5pt}
\setlength{\emergencystretch}{2em}
\numberwithin{equation}{section}
\newtheorem{theorem}{Theorem}[section]
\newaliascnt{proposition}{theorem}
\newtheorem{proposition}[proposition]{Proposition}
\aliascntresetthe{proposition}
\crefname{proposition}{Proposition}{Propositions}
\Crefname{proposition}{Proposition}{Propositions}
\newaliascnt{lemma}{theorem}
\newtheorem{lemma}[lemma]{Lemma}
\aliascntresetthe{lemma}
\crefname{lemma}{Lemma}{Lemmas}
\Crefname{lemma}{Lemma}{Lemmas}
\newaliascnt{corollary}{theorem}
\newtheorem{corollary}[corollary]{Corollary}
\aliascntresetthe{corollary}
\crefname{corollary}{Corollary}{Corollaries}
\Crefname{corollary}{Corollary}{Corollaries}
\theoremstyle{definition}
\newaliascnt{definition}{theorem}
\newtheorem{definition}[definition]{Definition}
\aliascntresetthe{definition}
\crefname{definition}{Definition}{Definitions}
\Crefname{definition}{Definition}{Definitions}
\newaliascnt{example}{theorem}
\newtheorem{example}[example]{Example}
\aliascntresetthe{example}
\crefname{example}{Example}{Examples}
\Crefname{example}{Example}{Examples}
\newaliascnt{problem}{theorem}
\newtheorem{problem}[problem]{Problem}
\aliascntresetthe{problem}
\crefname{problem}{Problem}{Problems}
\Crefname{problem}{Problem}{Problems}
\theoremstyle{remark}
\newaliascnt{remark}{theorem}
\newtheorem{remark}[remark]{Remark}
\aliascntresetthe{remark}
\crefname{remark}{Remark}{Remarks}
\Crefname{remark}{Remark}{Remarks}
\theoremstyle{plain}
\newtheorem{maintheorem}{Theorem}
\renewcommand{\themaintheorem}{\Alph{maintheorem}}
\crefname{maintheorem}{Theorem}{Theorems}
\Crefname{maintheorem}{Theorem}{Theorems}
\newcommand{\E}{\mathbf{Eff}}
\newcommand{\N}{\mathbb N}
\newcommand{\TV}{\mathbb T}
\newcommand{\Pow}{\mathcal P}
\newcommand{\OO}{\mathcal O}
\newcommand{\Pt}{\operatorname{Pt}_{\Sigma}}
\newcommand{\Spec}{\operatorname{Spec}}
\newcommand{\im}{\operatorname{im}}
\newcommand{\id}{\operatorname{id}}
\newcommand{\ev}{\operatorname{ev}}
\newcommand{\Sub}{\operatorname{Sub}}
\newcommand{\Hom}{\operatorname{Hom}}
\newcommand{\Adm}{\operatorname{Adm}}
\newcommand{\Par}{\operatorname{Par}}
\newcommand{\dom}{\operatorname{dom}}
\newcommand{\onto}{\twoheadrightarrow}
\newcommand{\incl}{\hookrightarrow}
\newcommand{\pmap}{\rightharpoonup}
\newcommand{\Cw}{\mathsf C_{\omega}}
\newcommand{\Mark}{\mathbf{Mark}_{\mathrm{mod}}}
\newcommand{\JoinMark}{\mathbf{Mark}_{\omega}}
\newcommand{\SSn}{\mathcal S}
\newcommand{\PPn}{\mathcal P}
\newcommand{\DDn}{\mathcal D}
\newcommand{\negT}{\neg_{\TV}}
\newcommand{\joinT}{\bigvee^{\TV}}
\newcommand{\meetT}{\bigwedge^{\TV}}
\newcommand{\jsem}{j^{\mathrm{sem}}}
\newcommand{\ObsCat}{\mathbf{Obs}_{\Sigma}}
\newcommand{\Profile}{\operatorname{Prof}}
\newcommand{\Tot}{\operatorname{Tot}}
\newcommand{\Icite}[1]{\cite[\textnormal{#1}]{sheridan}}
\newcolumntype{L}[1]{>{\raggedright\arraybackslash}p{#1}}
\setcounter{tocdepth}{1}
\title{\textcolor{red}{UNDER CONSTRUCTION} Effective Observation Hierarchies\\[4pt]\large A Synthetic Approach to Descriptive Set Theory}
\author{Conor Sheridan}
\date{November 2026}

\begin{document}
\maketitle

\begin{abstract}
We develop a finite descriptive hierarchy from the positive observation
classifier and observation duality of the effective topos. Interpreted
presentations first give a regular-cover descent theorem for predicates.
Countable join then defines a reflection on modest markings of the stable
truth object. Alternating this reflection with complement constructs
positive and negative observation levels, their predicate objects, and
their numerical quantifier operations. Characters of the inherited
positive observation algebras, interpreted in stable truth, reconstruct
the corresponding graded effective-point objects. The first level is the
effective-point construction of the preceding paper. We prove functoriality,
semantic product reconstruction, and exact admissibility formulas for
products and partial substitution. Termination and totality give intrinsic
separation examples. A subsequent comparison with the inherited numerical
presentation identifies the levels on the natural numbers with the finite
arithmetical hierarchy and proves strictness at every finite level.
The construction uses images, exponentials, adjunctions, and spectra;
numerical universality enters only the comparison and diagonal argument.
We distinguish the resulting graded reconstruction from the additional
endofunctor comparisons required to identify it with a standard limit jump.
\end{abstract}

\tableofcontents

\section{From positive observations to descriptive hierarchies}
\label{sec:introduction}

The preceding paper \cite{sheridan}, hereafter Paper~I, constructs semantic
domains and their effective points from a positive observation classifier
$i:\Sigma\incl\Omega$ in $\E$. Its two-spectrum theorem recovers both a
searchable semantic domain and its effective points from one positive
observation algebra. The present paper develops the observation side of
that construction beyond positivity. An observation may fail to be
positive while still admitting a structured presentation obtained from
positive observations by complement and countable operations. We construct
the objects which classify these observations and then reconstruct the
points that they admit.

The separation between semantic truth and effective observation remains
essential. The ordinary Boolean object $2$ classifies actual decisions in
$\E$. Paper~I proves that the computation object $P$ has no nonconstant
map to $2$. A descriptive construction must therefore distinguish a new
way of observing a Boolean value from an ordinary Boolean decision. We
make that distinction by interpreting observation algebras in the stable
truth object $\TV=\Omega_{\neg\neg}$ and retaining, at each level, a
specified modest marking of $\TV$.

Synthetic descriptive set theory studies descriptive classes through
modified observation objects and categorical operations. Pauly and
de~Brecht develop this perspective for represented spaces, including the
role of admissibility in recovering maps from inverse images of
observations \cite{paulydebrecht}. Our starting data are instead the
observation algebras and effective-point comparison already constructed in
Paper~I. We use the represented-space theory as a comparison, after the
internal construction has been made.

\begin{maintheorem}[Completion and the finite hierarchy]
\label{main:hierarchy}
The positive classifier embeds in $\TV$. On modest subobjects of $\TV$,
countable join defines a reflection into the countably join-closed
markings. Starting at $A_1=\Sigma$, complement and this reflection give
markings $A_n,B_n\incl\TV$, $n\geq1$, with
\[
 B_n=\negT(A_n),\qquad A_{n+1}=\Cw(B_n),\qquad
 A_n\incl A_{n+1},\quad B_n\incl A_{n+1}.
\]
For each fixed finite level, the associated predicate objects are natural
under substitution. The positive levels have finite meets and countable
joins; the negative levels have finite joins and countable meets.
The markings themselves form a single family indexed by the natural
numbers object.
\end{maintheorem}

\begin{maintheorem}[Graded reconstruction]
\label{main:spectra}
For a searchable domain $D$ with positive observation algebra $\Gamma_D$,
there are objects
\[
 \overline D=\Spec_{\TV}(\Gamma_D),\qquad
 D^{[n]}=\overline D\times_{\TV^{\Gamma_D}}A_n^{\Gamma_D}.
\]
The object $D^{[n]}$ represents level-$n$ observation profiles whose
interpretations are characters. The first level is canonically $\Pt D$.
These constructions are functorial for the effective Scott maps of
Paper~I, and the level inclusions are natural. Semantic spectra preserve
binary products. The corresponding graded product comparison is a mono
whose domain is described by an explicit universal admissibility formula.
Semantic maps likewise have maximal graded restrictions, with an exact
formula for their sequential domains.
\end{maintheorem}

\begin{maintheorem}[Separation and numerical comparison]
\label{main:separation}
Termination has a complementary level-$2$ observation, and it defines a
map $P\to2^{[2]}$ which does not factor through $2^{[1]}\cong2$.
Totality is classified at the second negative level and is not positive.
On $\N$, the supports of level-$n$ predicates are exactly the usual
$\Sigma^0_n$ and $\Pi^0_n$ numerical predicates. There is a universal
numerical predicate for each positive level, and its complemented diagonal
proves strictness of the positive hierarchy at every finite level.
\end{maintheorem}

The proofs of the first two theorems precede the numerical comparison.
The latter uses the program cover already derived in Paper~I; it is not
used to define the hierarchy or its spectra. The exact product and
substitution results also specify the remaining comparison problem: a
graded reconstruction is not yet an endofunctor on a category closed under
that reconstruction. We retain that distinction in \cref{sec:comparison}.

\section{Inherited structure and stable truth}
\label{sec:truth}

\subsection{The ground model and imported results}
\label{subsec:imports}

We work throughout in the fixed effective topos of Paper~I, with the
realizability interpretation of \cite{hyland,vanoosten}. Its natural
numbers object is $\N$, its positive classifier is $i:\Sigma\incl\Omega$,
and its ordinary Boolean object is $2=1+1$. Images are regular images.
Predicates and existential elimination therefore have their exact-category
meaning; an image factorization does not choose a section of its cover.

We use the finite-hit cover $e:2^{\N}\onto\Sigma$, numerical
$\Sigma$-overtness, and the finite positive operations from
\Icite{Finite-hit classification and Relative foundation}. We also use
\Icite{Function names and subobjects, Modest-target lifting, and Finite-hit
name presentations}: $\Sigma$ is modest, modest objects are closed under
finite products, exponentials between modest objects, and subobjects,
and a map from a partitioned assembly lifts through a regular epimorphism
whose domain is modest. This last lifting theorem uses unique denotation
of a supplied name and does not use a selection of arbitrary fibres.

The searchable presentations, positive trace algebras, and observation
duality are those of \Icite{Searchable approximation presentations and
Two-spectrum reconstruction}. The intrinsic computation object is denoted
by $P$, and $H:P\times\N\to\Sigma$ is its termination observation.
We use \Icite{Intrinsic Rice and halting, Intrinsic finite information,
and the nonpositivity corollary} for the separation examples.
\Cref{app:imports} records the original source labels of these results.

We reserve $\TV$ for stable truth, leaving the notation for Baire space
in Paper~I available for later comparisons. A base carrier of a searchable
presentation is written $B$ or $C$; it is not the stable truth object.

\subsection{The Boolean algebra of stable truth}

\begin{definition}[Semantic truth]
\label{def:semantic-truth}
Let $j=\neg\neg:\Omega\to\Omega$. Write
\[
 \TV=\{p:\Omega\mid j(p)=p\},\qquad k:\TV\incl\Omega.
\]
For a family $p:A\to\TV$, define
\begin{equation}\label{eq:stable-operations}
 \negT p=\neg p,\qquad
 \meetT_{a:A}p(a)=\forall a:A\;p(a),\qquad
 \joinT_{a:A}p(a)=j\bigl(\exists a:A\;p(a)\bigr).
\end{equation}
\end{definition}

\begin{lemma}[Stable truth operations]
\label{lem:stable-boolean}
The operations in \eqref{eq:stable-operations} make $\TV$ an internally
complete Boolean algebra. Finite meets distribute over its indexed joins.
Its indexed operations satisfy substitution and Fubini equations.
The map $k$ preserves meets, but its joins are the $j$-closures of ambient
joins.
\end{lemma}
\begin{proof}
A universal quantification of stable predicates is stable: from
$\neg\neg\forall a\,p(a)$ one obtains $\neg\neg p(a)$ for each $a$, hence
$p(a)$. Negations and double negations are stable. The displayed meet is
the greatest stable lower bound; the displayed join is the least stable
upper bound, since $\exists a\,p(a)\leq q=j(q)$ implies
$j(\exists a\,p(a))\leq q$. Complement satisfies
$p\wedge\neg p=\bot$ and $j(p\vee\neg p)=\top$, and is involutive on
stable truth. The nucleus $j$ preserves finite meets. Therefore, for stable
$q$, $q\wedge j(r)=j(q\wedge r)$, which proves distributivity over
indexed joins. Substitution follows from the corresponding subobject
quantifiers, and interchange of existential quantifiers followed by
idempotence of $j$ proves the join Fubini equation. The meet equation is
the interchange of universal quantifiers.
\end{proof}

This Boolean algebra lives in the constructive topos $\E$. Its equation
$p\vee^{\TV}\negT p=\top$ does not assert the ambient proposition
$p\vee\neg p$. This distinction will account for the changed Boolean
target in \cref{sec:examples}.

\begin{proposition}[Positive truth is stable]
\label{prop:positive-stability}
There is a unique mono $\alpha:\Sigma\incl\TV$ with $k\alpha=i$.
It preserves finite meets and all the overt joins defined in $\Sigma$.
\end{proposition}
\begin{proof}
An assembly is separated for double negation; equivalently its equality
predicate is stable. This can be seen directly in the assembly
presentation: in a context carrying names of $x$ and $y$, nonempty
double-negated equality forces equality of their underlying elements,
and the supplied name of $x$ then realizes their assembly equality.
This is uniform in the given context. The categorical characterization
is recalled in \cite[Section~3.2]{hylandvariations}.

The object $\Sigma$ is modest and hence an assembly. For $p:\Sigma$,
$i(p)$ is equivalent to $p=\top_\Sigma$. Thus $i(p)$ is stable, giving
the factorization, and monicity of $i$ makes $\alpha$ monic. For a
$\Sigma$-overt $A$ and $u:A\to\Sigma$, the ordinary existential
$\exists a\,i(u(a))$ already has its $\Sigma$-classifier
$\bigvee_A^\Sigma u$. It is stable by the first part. Taking its
$j$-closure changes nothing, proving preservation of that join.
Finite meets are preserved by both inclusions.
\end{proof}

The interpretation is the commuting triangle
\begin{equation}\label{diag:truth-factorization}
\begin{tikzcd}[column sep=large]
 \Sigma\ar[r,hook,"\alpha"]\ar[dr,hook,"i"']&
 \TV\ar[d,hook,"k"]\\
 &\Omega.
\end{tikzcd}
\end{equation}
Equalities between positive operations can consequently be checked in
$\TV$ and reflected through $\alpha$.

\section{Interpreted presentations and regular descent}
\label{sec:descent}

\begin{definition}[Markings and presentations]
\label{def:interpreted-presentation}
A \emph{marking} of an object $H$ is a mono $m:M\incl H$.
An \emph{interpreted presentation} of this marking is a factorization
$\tau=mq:T\onto M\incl H$ with $q$ regular epic.
A predicate $p:X\to H$ is $M$-classified when it factors through $m$.
An actual presentation of $p$ is a map $t:X\to T$ with $\tau t=p$.
\end{definition}

\begin{theorem}[Descent of presented predicates]
\label{thm:predicate-descent}
For an interpreted presentation as above and $p:X\to H$, the following
are equivalent:
\begin{enumerate}
 \item $p$ is $M$-classified;
 \item there are a regular epi $r:Y\onto X$ and a map $t:Y\to T$
       such that $pr=\tau t$.
\end{enumerate}
The equivalence holds in arbitrary parameter contexts and is stable under
pullback. The $M$-classified predicates form the least class containing
the actually presented predicates and closed under descent along regular
covers of their domains.
\end{theorem}
\begin{proof}
For a factorization $p=m\bar p$, use the canonical pullback
\begin{equation}\label{diag:presentation-cover}
\begin{tikzcd}[column sep=large]
 Y=X\times_M T\ar[r,"t"]\ar[d,two heads,"r"']&
 T\ar[d,two heads,"q"]\\
 X\ar[r,"\bar p"']&M\ar[r,hook,"m"']&H.
\end{tikzcd}
\end{equation}
Regular epimorphisms are stable under pullback, so $r$ has the required
property. Conversely let $U=X\times_H M\incl X$. The equality
$pr=\tau t$ makes $r$ factor through $U$. Its image is all of $X$,
so the mono $U\incl X$ is an isomorphism. This gives the unique
factorization through $M$.

Replace $X$ by $Z\times X$ to retain a parameter $Z$. Both directions
use pullbacks and stable images, so reindexing in $Z$ preserves them.
Every $M$-classified predicate is locally presented by the displayed
cover. The $M$-classified predicates themselves descend, by the converse
argument. These two observations prove the leastness assertion.
\end{proof}

\begin{proposition}[Presentation objects and predicate objects]
\label{prop:presentation-comparison}
There is a canonical mono
\begin{equation}\label{eq:presentation-comparison}
 \im(\tau^X)\incl M^X\incl H^X.
\end{equation}
The first mono is an isomorphism precisely when $q^X$ is regular epic.
Membership in its left-hand side is local lifting of a whole family of
predicates over its parameter object. Membership in $M^X$ permits the
domain-cover descent of \cref{thm:predicate-descent}.
\end{proposition}
\begin{proof}
Exponentiation preserves the mono $m$, and $\tau^X=m^Xq^X$.
Its image is therefore $m^X\im(q^X)$, giving the mono and the stated
criterion. For a parameter $Z\to H^X$, the image criterion for
$\tau^X$ asks for a regular cover of $Z$ on which the arrow lifts to
$T^X$. In contrast, factoring its transpose $Z\times X\to H$ through
$M$ is equivalent to lifting after a regular cover of $Z\times X$.
These are the two asserted local conditions.
\end{proof}

\begin{corollary}[Descent completion]
\label{cor:descent-completion}
The marking $M\incl H$ depends only on the image of the interpretation
$\tau$. Replacing $T$ by any other regular presentation of the same
image leaves all completed predicate objects $M^X$ and their substitution
maps canonically unchanged.
\end{corollary}
\begin{proof}
Images are unique up to the unique isomorphism over $H$. Exponentiation
and precomposition transport that isomorphism and preserve its uniqueness.
\end{proof}

\section{Countable-join completion}
\label{sec:completion}

\subsection{The covers used by countable operations}

\begin{lemma}[Modest images in assemblies]
\label{lem:modest-image}
The regular image of a map from a modest object to an assembly is modest.
In particular, the image of a modest object in $\TV$ is modest.
\end{lemma}
\begin{proof}
Write the map as a tracked map of assemblies $f:D\to S$. On the
set-theoretic image of $f$, declare a name $a$ to realize $y$ when $a$
realizes a $d\in D$ with $f(d)=y$. The relation is inhabited at every
image value. Since $D$ is modest, one name determines at most one $d$,
and hence at most one $y$. Thus this is a modest assembly $M$.
The map $D\to M$ is tracked by the identity on names. A supplied
$M$-name already supplies a name of a preimage, so the map is regular
epic. The inclusion $M\to S$ is tracked by the tracker of $f$ and is
monic. This is the regular-image factorization, also in $\E$ under the
regular inclusion of assemblies. Finally $\TV$ is a double-negation
sheaf, hence a separated object and an assembly. Concretely it is the
codiscrete Boolean assembly under the standard identification of the
double-negation sheaves of $\E$ with sets.
\end{proof}

\begin{lemma}[Countable lifting]
\label{lem:countable-cover-lifting}
If $q:T\onto M$ is a regular epi with $T$ modest, then
$q^{\N}:T^{\N}\onto M^{\N}$ is regular epic.
\end{lemma}
\begin{proof}
Given $z:Z\to M^{\N}$, take a canonical partitioned cover
$R\onto Z$. The transpose of the induced map $R\to M^{\N}$ has
partitioned domain $R\times\N$. Modest-target lifting from Paper~I
gives $R\times\N\to T$ above it. Transposition gives a lift
$R\to T^{\N}$. In particular, applying this construction to
$Z=M^{\N}$ and $z=\id$ proves that the image of $q^{\N}$ is all
of $M^{\N}$. No lifts were selected as a function of arbitrary elements
of $Z$.
\end{proof}

\subsection{A reflection on marked truth}

Let $\Mark(\TV)$ be the poset of modest subobjects of $\TV$, with
inclusions over $\TV$ as arrows. It is an ordinary poset of subobjects;
no assertion that it is a small internal category is needed below.
Let $\JoinMark(\TV)$ be its full subposet of markings through which
countable semantic join restricts.

\begin{definition}[Countable completion]
\label{def:countable-completion}
For $m:M\incl\TV$ define $\Cw(M)$ by
\begin{equation}\label{eq:countable-image}
 M^{\N}\xrightarrow{q_M}\Cw(M)\xrightarrow{c_M}\TV,
 \qquad c_Mq_M=\joinT_{\N}m^{\N},
\end{equation}
where $q_M$ is regular epic and $c_M$ monic.
\end{definition}

\begin{theorem}[Countable completion reflection]
\label{thm:countable-completion-reflection}
The construction $\Cw$ defines a left adjoint to
$\JoinMark(\TV)\incl\Mark(\TV)$. Its unit is the inclusion obtained
from constant sequences. Consequently it is monotone and idempotent,
with all the comparison equations determined by the inclusions in $\TV$.
\end{theorem}
\begin{proof}
The object $M^{\N}$ is modest, so \cref{lem:modest-image} makes
$\Cw(M)$ modest. Constant sequences and inhabitation of $\N$ give
$M\incl\Cw(M)$.

To construct countable join on $\Cw(M)$, exponentiate $q_M$ by $\N$.
The resulting map is epic by \cref{lem:countable-cover-lifting}. On its
domain $(M^{\N})^{\N}\cong M^{\N\times\N}$, flatten a double sequence
using the fixed numerical pairing isomorphism and apply $q_M$. Its
interpretation is
\[
 \joinT_i\joinT_j m(u_{ij})=\joinT_{(i,j)}m(u_{ij}).
\]
Hence two inputs with the same image in $\Cw(M)^{\N}$ have the same
output after $c_M$, and therefore the same output. Effectiveness of the
kernel pair descends the map to $\Cw(M)^{\N}\to\Cw(M)$.
Its interpretation is the desired join, and monicity gives uniqueness.

If $M\incl L\incl\TV$ and $L$ is countably join-closed, then the
interpretation $M^{\N}\to\TV$ factors through $L$. Its image
$\Cw(M)$ therefore factors through $L$. Conversely
$\Cw(M)\incl L$ implies $M\incl L$ by the unit. This is the
adjunction. The universal property gives monotonicity and the inverse
inclusions for $\Cw\Cw(M)=\Cw(M)$. All equations between its maps
hold after the monomorphism into $\TV$, so hold before it.
\end{proof}

The universal factorization is displayed by
\begin{equation}\label{diag:completion-universal}
\begin{tikzcd}[column sep=large]
 M^{\N}\ar[r,two heads,"q_M"]\ar[dr]&
 \Cw(M)\ar[d,hook,dashed]\ar[dr,hook,"c_M"]&\\
 &L\ar[r,hook]&\TV.
\end{tikzcd}
\end{equation}
The dashed arrow is determined by the marking, not by a chosen
normal-form function.

\begin{lemma}[Finite meets through completion]
\label{lem:completion-meets}
If a modest marking $M$ contains $\bot,\top$ and is closed under finite
meets in $\TV$, then $\Cw(M)$ has these properties and is closed under
finite joins.
\end{lemma}
\begin{proof}
Constants belong by the unit. Finite joins follow from countable join,
using $\bot$ for unused entries. On two presented inputs, define their
meet by the doubly indexed family $u_i\wedge v_j$ in $M$. Its
interpretation is
\[
 \left(\joinT_i u_i\right)\wedge
 \left(\joinT_j v_j\right)=\joinT_{i,j}(u_i\wedge v_j).
\]
The product of the presentation covers is epic. The displayed equality
and monicity of $c_M$ descend the construction through its kernel pair.
The nullary meet is the already included $\top$.
\end{proof}

\begin{remark}[Completion and descriptive successor]
\label{rem:successor-not-reflection}
The involution $\negT:\TV\to\TV$ induces an order automorphism of
the poset of markings by taking direct images. The descriptive successor
will be $M\mapsto\Cw(\negT M)$. Idempotence of $\Cw$ does not imply
idempotence of this composite. In particular, completion under one
operation and alternation of operations are different constructions.
\end{remark}

\section{The finite hierarchy and its predicate objects}
\label{sec:hierarchy}

\subsection{An indexed construction of the markings}

\begin{definition}[Observation levels]
\label{def:observation-hierarchy}
Regard $\alpha:\Sigma\incl\TV$ as the marking $A_1$. For $n\geq1$
define $B_n\incl\TV$ by transporting $A_n$ along $\negT$, and define
$A_{n+1}$ as the image of countable join on $B_n$. Denote the markings
by $a_n:A_n\incl\TV$ and $b_n:B_n\incl\TV$, and the successor cover
by
\begin{equation}\label{eq:successor-cover}
 q_{n+1}:B_n^{\N}\onto A_{n+1},\qquad
 a_{n+1}q_{n+1}=\joinT_{\N}b_n^{\N}.
\end{equation}
\end{definition}

\begin{proposition}[Internal indexing]
\label{prop:hierarchy-indexing}
The markings of \cref{def:observation-hierarchy} are the fibres of
subobjects $\mathcal A,\mathcal B\incl\N\times\TV$, with the fibre
at $r$ denoted $A_{r+1},B_{r+1}$. Their image presentations are obtained
from one construction in the slice over $\N$.
\end{proposition}
\begin{proof}
Put $\Pow(\TV)=\Omega^{\TV}$. The universal membership predicate
defines an arrow $c:\Pow(\TV)\to\Pow(\TV)$ whose value on a marking
$U$ is the image of countable join of sequences in $U$. Explicitly,
\begin{equation}\label{eq:powerobject-completion}
 c(U)(p)\;\longleftrightarrow\;
 \exists s:\TV^{\N}\,
 \bigl[(\forall k\;U(s_k))\wedge p=\joinT_k s_k\bigr].
\end{equation}
This formula specifies a regular image of a universal family and hence
an arrow of the topos. The complement action is the arrow
$v(U)(p)=U(\negT p)$. Numerical recursion constructs the unique
$F:\N\to\Pow(\TV)$ with $F(0)=A_1$ and $F(r+1)=c(v(F(r)))$.
Pull back universal membership along $F\times\id$ to obtain
$\mathcal A$; use $vF$ for $\mathcal B$.

The dependent object of sequences in the fibres of $\mathcal B$ and
its interpretation by countable join give a single image factorization
over $\N$. Pullback along a numeral identifies it with
\eqref{eq:successor-cover}, by stability of images. This proves both
the indexing and compatibility of the presentations.
\end{proof}

The construction in \cref{prop:hierarchy-indexing} is an internal
family, not merely a meta-level list. The closure theorems below are
proved for each fixed finite level, with arbitrary object parameters.
They do not assert a uniform choice of presentations across levels, or
an internally quantified modestness principle for arbitrary markings.

\begin{lemma}[Modesty and the first inclusion]
\label{lem:hierarchy-base}
Each fixed finite $A_n$ and $B_n$ is modest. There is a canonical
inclusion $A_1\incl A_2$ over $\TV$.
\end{lemma}
\begin{proof}
Modesty starts with $\Sigma$. Complement transports it by an
isomorphism, exponentiation by $\N$ preserves it, and
\cref{lem:modest-image} applies to the successor image. This proves
the first assertion by induction on the finite level.

Let $\beta:2\to\Sigma$ be the Boolean positive classifier. Each
Boolean value lies in $B_1$: complement its Boolean value, apply
$\beta$, and then use the complement isomorphism $A_1\cong B_1$.
This defines $z:2^{\N}\to B_1^{\N}$. Its interpretation by countable
join is $\alpha e$, because $e$ is the finite-hit classifier and
$\alpha$ preserves the numerical positive join. Hence $\alpha e$
factors through $A_2$. Since $e$ is epic, its image contains the image
of $\alpha$, proving the inclusion. This uses the whole finite-hit
cover, not a section of it.
\end{proof}

\begin{theorem}[Hierarchy operations]
\label{thm:hierarchy-operations}
For each fixed $n\geq1$ there are inclusions
\[
 A_n\incl A_{n+1},\qquad B_n\incl B_{n+1},\qquad
 B_n\incl A_{n+1},\qquad A_n\incl B_{n+1}.
\]
The object $A_n$ contains the Boolean constants and is closed under
finite meets and countable semantic joins. The object $B_n$ contains
the Boolean constants and is closed under finite joins and countable
semantic meets. All displayed operations and inclusion equations are
morphisms and commute with arbitrary parameter substitution.
\end{theorem}
\begin{proof}
Monotonicity of complement on markings and of $\Cw$, applied
inductively to \cref{lem:hierarchy-base}, gives $A_n\incl A_{n+1}$.
Complement gives $B_n\incl B_{n+1}$. The unit of $\Cw$ gives
$B_n\incl A_{n+1}$, and complement gives $A_n\incl B_{n+1}$.

At the base, the asserted positive operations are those of $\Sigma$,
interpreted by \cref{prop:positive-stability}. If $A_n$ has finite
joins, its complement marking $B_n$ has finite meets.
\Cref{thm:countable-completion-reflection,lem:completion-meets}
therefore give the positive operations on $A_{n+1}$. Complement gives
the negative operations. This proves the induction. Each operation was
constructed by descent as an actual arrow, and its interpretation is
the corresponding $\TV$-operation. Equality after the marking proves
its laws and its substitution equations in every parameter context.
\end{proof}

\subsection{Predicate objects and quantifiers}

\begin{definition}[Positive, negative, and ambiguous predicates]
\label{def:predicate-levels}
For $X\in\E$ put
\[
 \SSn_n(X)=A_n^X,\qquad \PPn_n(X)=B_n^X,
 \qquad \Delta_n=A_n\times_{\TV}B_n,
 \qquad \DDn_n(X)=\Delta_n^X.
\]
All these objects are marked subobjects of $\TV^X$. A level-$n$
predicate is a map classified by the indicated marking.
\end{definition}

\begin{proposition}[Substitution and numerical quantification]
\label{prop:quantifier-calculus}
The predicate objects define contravariant functors in $X$. Their
finite operations are pointwise, and numerical quantification has the
types
\begin{align}
 \exists^{\TV}_{\N}&:\SSn_n(X\times\N)\to\SSn_n(X),&
 \forall^{\TV}_{\N}&:\PPn_n(X\times\N)\to\PPn_n(X),
 \label{eq:same-level-quantifiers}\\
 \exists^{\TV}_{\N}&:\PPn_n(X\times\N)\to\SSn_{n+1}(X),&
 \forall^{\TV}_{\N}&:\SSn_n(X\times\N)\to\PPn_{n+1}(X).
 \label{eq:cross-level-quantifiers}
\end{align}
These maps satisfy the Beck--Chevalley equations for substitution in
$X$, the corresponding Fubini equations, and the finite-meet Frobenius
equation for existential quantification whenever the displayed types
admit its two sides.
\end{proposition}
\begin{proof}
Precomposition defines substitution. Transpose a predicate on
$X\times\N$ to a family on $X$, and apply the operations of
\cref{thm:hierarchy-operations}. The first cross-level map is
$q_{n+1}$ after transposition; its dual gives the second.
The equations hold in $\TV$ by \cref{lem:stable-boolean}.
The relevant markings are monic, so those equations descend to the
indicated predicate objects.
\end{proof}

The existential in these formulas is $j\exists$, not the unrestricted
ambient existential. No identification of these two operations on the
higher levels is made. At the positive base and for a $\Sigma$-overt
index, \cref{prop:positive-stability} gives that identification.

\begin{proposition}[The ambiguous base]
\label{prop:ambiguous-base}
Each $\Delta_n$ is closed under the finite Boolean operations interpreted
in $\TV$, and $\Delta_1\cong2$ over $\TV$.
\end{proposition}
\begin{proof}
Both markings in the pullback are closed under finite meets and finite
joins. Complement exchanges them, so restricts to their intersection.
For $n=1$, a generalized element of $\Delta_1$ supplies positive
predicates $p,q$ whose interpreted values are complements. Their meet
and join in $\Sigma$ are $\bot$ and $\top$, since $\alpha$ preserves
and reflects those operations. Thus they form an actual complementary
pair of subobjects of the parameter object. Its coproduct decomposition
gives a unique Boolean decision. Conversely a Boolean decision supplies
that pair. The two constructions are inverse and natural in the
parameter object, proving the isomorphism.
\end{proof}

\begin{proposition}[Universal positive projection]
\label{prop:universal-classifier}
The map
\[
 u=\forall_{\N}i^{\N}:\Sigma^{\N}\longrightarrow\Omega
\]
has image $B_2\incl\TV\incl\Omega$. Consequently the regular-cover
descent completion of its presented predicates on $X$ is $B_2^X$.
\end{proposition}
\begin{proof}
Universal quantification of stable predicates is stable. In $\TV$,
De Morgan duality expresses $u$ as the complement of countable join
of the complementary $B_1$-values. The image of that join is $A_2$.
Transport of an image through the complement isomorphism therefore
gives $B_2$. Apply \cref{thm:predicate-descent} to this image
factorization. This does not identify the raw image of $u^X$ with
$B_2^X$ without the additional lifting property of
\cref{prop:presentation-comparison}.
\end{proof}

\section{Spectra of the inherited observation algebras}
\label{sec:spectra}

\subsection{Characters as an internally defined object}

Let $(B,r)$ be a searchable approximation presentation of Paper~I, with
$\rho(a,b)\in\Sigma$ its approximation relation. Its positive trace
object is
\begin{equation}\label{eq:trace-object}
 \Gamma_r=\left\{u:B\to\Sigma\ \middle|\
 u(b)=\bigvee_{a:B}^{\Sigma}u(a)\wedge\rho(a,b)\right\}.
\end{equation}
Write $D=D_r$, $\Gamma_D=\Gamma_r$, and $\delta_a(b)=\rho(a,b)$.
The positive basis identities proved in Paper~I are
\begin{align}
 \top&=\bigvee_{b:B}^{\Sigma}\delta_b,
 \label{eq:basis-top}\\
 \delta_a&=\bigvee_{c:B}^{\Sigma}\rho(a,c)\wedge\delta_c,
 \label{eq:basis-round}\\
 \delta_a\wedge\delta_b
 &=\bigvee_{c:B}^{\Sigma}\rho(a,c)\wedge\rho(b,c)\wedge\delta_c,
 \label{eq:basis-meet}\\
 u&=\bigvee_{b:B}^{\Sigma}u(b)\wedge\delta_b.
 \label{eq:basis-expansion}
\end{align}
Joins in these equations are the joins of the trace algebra, computed
pointwise. The base $B$ is overt; its order need not itself be positive.

\begin{definition}[Semantic characters]
\label{def:semantic-spectrum}
The object $\overline D=\Spec_{\TV}(\Gamma_D)$ is the subobject of
$\TV^{\Gamma_D}$ of maps $\chi$ satisfying the following equations:
\begin{align}
 \chi(\underline p)&=\alpha(p) &&(p:\Sigma),
 \label{eq:char-constants}\\
 \chi(u\wedge v)&=\chi(u)\wedge\chi(v) &&(u,v:\Gamma_D),
 \label{eq:char-meets}\\
 \chi(u)&=\joinT_{b:B}\alpha(u(b))\wedge\chi(\delta_b)
 &&(u:\Gamma_D).
 \label{eq:char-expansion}
\end{align}
Here $\underline p$ is the constant trace. The defining equations
quantify over the displayed objects, so define an actual subobject of
the exponential.
\end{definition}

\begin{proposition}[Character operations]
\label{prop:character-operations}
The characters of \cref{def:semantic-spectrum} are precisely the maps
preserving finite meets, $\Sigma$-constants, and all overt-indexed joins,
with these operations interpreted in $\TV$.
\end{proposition}
\begin{proof}
Preservation of constants supplies top, so
\eqref{eq:char-meets} supplies finite meets. If $A$ is overt and
$v:A\to\Gamma_D$, then
\begin{align*}
 \chi\left(\bigvee_{a:A}^{\Sigma}v_a\right)
 &=\joinT_{b:B}
   \alpha\left(\bigvee_{a:A}^{\Sigma}v_a(b)\right)
   \wedge\chi(\delta_b)\\
 &=\joinT_{a:A}\joinT_{b:B}
   \alpha(v_a(b))\wedge\chi(\delta_b)
 =\joinT_{a:A}\chi(v_a).
\end{align*}
The middle equality uses \cref{prop:positive-stability} and Fubini.
Conversely a map preserving the stated operations sends
\eqref{eq:basis-expansion} to \eqref{eq:char-expansion}; it also
satisfies the first two defining equations. Thus preservation of a
potentially large collection of indexed operations has been reduced
to the explicitly defined character object.
\end{proof}

\begin{theorem}[Stable basis reconstruction]
\label{thm:stable-basis}
The object $\overline D$ is canonically isomorphic to the object of
maps $I:B\to\TV$ satisfying
\begin{align}
 \joinT_b I(b)&=\top,
 \label{eq:stable-ideal-top}\\
 I(a)&=\joinT_c\alpha(\rho(a,c))\wedge I(c),
 \label{eq:stable-ideal-round}\\
 I(a)\wedge I(b)&=
 \joinT_c\alpha(\rho(a,c))\wedge\alpha(\rho(b,c))\wedge I(c).
 \label{eq:stable-ideal-meet}
\end{align}
The inverse maps are
\begin{equation}\label{eq:stable-basis-inverses}
 I_\chi(b)=\chi(\delta_b),\qquad
 \chi_I(u)=\joinT_b I(b)\wedge\alpha(u(b)).
\end{equation}
They are natural in arbitrary parameters.
\end{theorem}
\begin{proof}
Applying a character to
\cref{eq:basis-top,eq:basis-round,eq:basis-meet} gives the three
equations for $I_\chi$. Conversely let $I$ satisfy those equations.
Its interpretation on a constant trace is
$\alpha(p)\wedge\joinT_b I(b)=\alpha(p)$. For finite meets, expand
the two factors and apply \eqref{eq:stable-ideal-meet}:
\begin{align*}
 \chi_I(u)\wedge\chi_I(v)
 &=\joinT_{a,b}I(a)\wedge I(b)\wedge\alpha(u(a))\wedge\alpha(v(b))\\
 &=\joinT_c I(c)\wedge
   \left(\joinT_a\alpha(u(a)\wedge\rho(a,c))\right)
   \wedge
   \left(\joinT_b\alpha(v(b)\wedge\rho(b,c))\right)\\
 &=\joinT_c I(c)\wedge\alpha(u(c))\wedge\alpha(v(c))
 =\chi_I(u\wedge v).
\end{align*}
The last equality uses the rounding equation of the traces and
preservation of overt joins by $\alpha$. Further,
$\chi_I(\delta_a)=I(a)$ by \eqref{eq:stable-ideal-round}.
Consequently \eqref{eq:char-expansion} follows from the definition of
$\chi_I$. Thus $\chi_I$ is a character. Evaluation at $\delta_b$
proves $I_{\chi_I}=I$, and the expansion equation proves
$\chi_{I_\chi}=\chi$. All the expressions are arrows with their
parameters retained, proving naturality.
\end{proof}

The equations in this theorem are equations of stable truth. In
particular, \eqref{eq:stable-ideal-top} asserts stable inhabitation of
the basis predicate. It does not replace the ordinary inhabitation
condition in the definition of the semantic domain $D_r$ of Paper~I.
This is the precise reason to distinguish $\overline D$ from $D_r$.

\subsection{Marked spectra and the recovery theorem}

\begin{definition}[Graded effective points]
\label{def:descriptive-spectrum}
For $n\geq1$, define $D^{[n]}$ by the pullback
\begin{equation}\label{diag:graded-spectrum}
\begin{tikzcd}[column sep=large,row sep=large]
 D^{[n]}\ar[r,hook,"\lambda_{D,n}"]\ar[d,hook,"\jsem_{D,n}"']&
 A_n^{\Gamma_D}\ar[d,hook,"a_n^{\Gamma_D}"]\\
 \overline D\ar[r,hook,"\varepsilon_D"']&\TV^{\Gamma_D}.
\end{tikzcd}
\end{equation}
Write $\Profile_n(X,D)$ for the subobject of $A_n^{X\times\Gamma_D}$
whose interpreted transpose is a family of semantic characters.
\end{definition}

\begin{theorem}[Representation of observation profiles]
\label{thm:profile-representation}
There is a canonical isomorphism of objects
\begin{equation}\label{eq:profile-object}
 (D^{[n]})^X\cong\Profile_n(X,D).
\end{equation}
On arrows it identifies maps $X\to D^{[n]}$ with maps
$h:\Gamma_D\to A_n^X$ whose interpreted transpose preserves the
observation operations. The isomorphism is natural in $X$.
\end{theorem}
\begin{proof}
Exponentiate \eqref{diag:graded-spectrum} by $X$. Exponentiation is a
right adjoint and preserves the pullback. The exponential isomorphisms
$(A_n^{\Gamma_D})^X\cong A_n^{X\times\Gamma_D}$ and their semantic
counterparts identify the resulting pullback with the defining
subobject $\Profile_n(X,D)$. Explicitly, a map to $D^{[n]}$ gives its
$A_n$-valued evaluation profile. In the other direction, the character
condition supplies the map into $\overline D$, and the pullback gives
the unique lift. Evaluation and transposition are inverse, so both
composites are identities. Precomposition in $X$ commutes with every
step.
\end{proof}

\begin{theorem}[Recovery of the original effective points]
\label{thm:base-recovery}
There is a natural isomorphism $D^{[1]}\cong\Pt D$.
\end{theorem}
\begin{proof}
A profile at level $1$ is a map $\Gamma_D\to\Sigma$ whose
interpretation through $\alpha$ satisfies the character equations.
The operations on its positive values are preserved by $\alpha$.
Monicity therefore reflects these equations, so the profile is exactly
a $\Sigma$-character. Conversely every such character gives a level-$1$
profile. The two constructions leave the profile unchanged and are
inverse. Apply the two-spectrum theorem of Paper~I,
$\Spec_\Sigma(\Gamma_D)\cong\Pt D$. Its naturality and the unchanged
profiles prove the stated naturality.
\end{proof}

\begin{proposition}[The computational-carrier obstruction]
\label{prop:computational-carrier-obstruction}
Let $T\to\Omega$ be any interpretation and form
$E_D=D\times_{\Omega^{\Gamma_D}}T^{\Gamma_D}$ using the original
semantic profile of $D$. Every map $X\to E_2$ induces an ordinary
decision $X\to2$. In particular such a lift cannot interpret
termination on $P$ as its projected Boolean value.
\end{proposition}
\begin{proof}
The pullback has a projection $E_2\to2$. Composition gives the
decision. If it classified termination, its values at the
nowhere-defined function and a total constant function would differ,
contradicting intrinsic Rice in Paper~I.
\end{proof}

Our pullback \eqref{diag:graded-spectrum} projects to
$\Spec_{\TV}(\Gamma_D)$. Inclusion $\TV\incl\Omega$ does not in
general preserve the joins in the character equations, so it does not
give a map from this spectrum to $\Spec_\Omega(\Gamma_D)$.
The obstruction proposition thus identifies a condition that the
changed semantic interpretation genuinely avoids.

\section{Maps, products, and admissible substitution}
\label{sec:maps}

\subsection{Functoriality and independence of presentation}

Let $\ObsCat$ denote the category of searchable domains with their
positive observation algebras and the effective Scott maps of Paper~I.
An arrow $f:D\to E$ has an observation transpose
$f^*:\Gamma_E\to\Gamma_D$ preserving the positive operations.

\begin{theorem}[Graded functoriality]
\label{thm:graded-functoriality}
Precomposition of characters defines functors
\[
 D\longmapsto\overline D,\qquad
 D\longmapsto D^{[n]}\quad:\ObsCat\longrightarrow\E.
\]
The inclusions $\iota_{D,n}:D^{[n]}\incl D^{[n+1]}$ form natural
transformations. All these constructions act jointly on parameterized
families of observation transposes. Isomorphic positive observation
algebras give canonically isomorphic graded spectra.
\end{theorem}
\begin{proof}
Send $\chi$ to $\chi f^*$. Preservation of the positive operations by
$f^*$ and \cref{prop:character-operations} make this a semantic
character. If its input profile is $A_n$-valued, precomposition is still
$A_n$-valued, so the defining pullbacks induce $f^{[n]}$.
For $D\xrightarrow fE\xrightarrow gF$,
$(gf)^*=f^*g^*$ and hence
$(\chi f^*)g^*=\chi(gf)^*$. This proves composition; the identity
equation is immediate. The same formulas are maps of the relevant
exponentials, so the action is joint in all parameters for which the
observation transpose is defined.

The mono $A_n\incl A_{n+1}$ induces $\iota_{D,n}$ by pullback.
The naturality square
\begin{equation}\label{diag:level-naturality}
\begin{tikzcd}[column sep=large]
 D^{[n]}\ar[r,"f^{[n]}"]\ar[d,hook,"\iota_{D,n}"']&
 E^{[n]}\ar[d,hook,"\iota_{E,n}"]\\
 D^{[n+1]}\ar[r,"f^{[n+1]}"']&E^{[n+1]}
\end{tikzcd}
\end{equation}
commutes because both profiles are $\chi f^*$. Coherence for multiple
level inclusions follows by the same monicity argument. Finally an
isomorphism of trace algebras and its inverse induce inverse maps on
characters and on each pullback.
\end{proof}

\subsection{Semantic products and their exact graded domain}

For product presentations, write $u\otimes v$ for the trace
$(b,c)\mapsto u(b)\wedge v(c)$. The product presentation of Paper~I
has base $B\times C$ and approximation relation $\rho\times\sigma$.
Its coordinate traces are $\delta_b\otimes\delta_c$.

\begin{theorem}[Semantic product reconstruction]
\label{thm:semantic-products}
There is a canonical isomorphism
\begin{equation}\label{eq:semantic-product}
 \overline{D\times E}\cong\overline D\times\overline E.
\end{equation}
For characters $\chi,\psi$, its inverse has profile
\begin{equation}\label{eq:product-profile}
 \Phi_{\chi,\psi}(w)=
 \joinT_{b:B,c:C}
 \chi(\delta_b)\wedge\psi(\delta_c)\wedge\alpha(w(b,c)).
\end{equation}
\end{theorem}
\begin{proof}
Restrict a product character along $u\mapsto u\otimes\top$ and
$v\mapsto\top\otimes v$ to obtain its two marginals. Conversely,
let $I(b)=\chi(\delta_b)$ and $J(c)=\psi(\delta_c)$. The function
$K(b,c)=I(b)\wedge J(c)$ satisfies the three stable basis equations
for the product relation. For the top equation, distribute the two
joins. For rounding, distribute the two rounding equations. For the
meet equation, distribute the two equations
\eqref{eq:stable-ideal-meet}. Thus \cref{thm:stable-basis} constructs
the product character, with profile \eqref{eq:product-profile}.

Its first marginal is $\chi$, since $\joinT_cJ(c)=\top$, and the
second is $\psi$. Conversely a product character $\theta$ satisfies
\[
 \theta(\delta_b\otimes\delta_c)
 =\theta(\delta_b\otimes\top)\wedge
  \theta(\top\otimes\delta_c).
\]
Thus it agrees on every coordinate with the character reconstructed
from its marginals. The expansion equation gives equality on every
trace. These are inverse morphisms, uniformly in parameters.
\end{proof}

\begin{theorem}[Exact graded product comparison]
\label{thm:graded-product-domain}
The canonical comparison
\[
 c_{D,E,n}:(D\times E)^{[n]}\longrightarrow D^{[n]}\times E^{[n]}
\]
is monic. Under \eqref{eq:semantic-product}, its image is the subobject
$K_n(D,E)$ defined by
\begin{equation}\label{eq:product-admission}
 (\chi,\psi)\in K_n(D,E)quad\longleftrightarrow\quad
 \forall w:\Gamma_{D\times E}\;
       \Phi_{\chi,\psi}(w)\in A_n.
\end{equation}
In particular $(D\times E)^{[n]}\cong K_n(D,E)$, naturally in both
variables. At level $1$ the comparison is an isomorphism.
\end{theorem}
\begin{proof}
Projections induce the comparison by \cref{thm:graded-functoriality}.
After inclusion into the semantic spectra it is the isomorphism
\eqref{eq:semantic-product}, so it is monic. Its profile belongs to
$A_n^{\Gamma_{D\times E}}$ exactly when all its values do. This is
\eqref{eq:product-admission}. Monicity of $a_n$ gives the unique
joint factorization, and the pullback
\eqref{diag:graded-spectrum} supplies the inverse on that subobject.
The construction is natural because restriction of characters commutes
with projections and with their reconstruction. At the base level use
\cref{thm:base-recovery} and the effective product theorem of Paper~I.
\end{proof}

The quantifier over $w$ in \eqref{eq:product-admission} is an ordinary
subobject condition. It is not asserted to be a positive predicate.
The theorem gives the full comparison map and its maximal domain of
inverse reconstruction. Countable-join closure alone does not remove
that condition: the joins in \eqref{eq:product-profile} range over
the overt bases $B$ and $C$, which need not be $\N$.

Equivalently, the admissibility subobject is the explicit pullback
\begin{equation}\label{diag:product-admission}
\begin{tikzcd}[column sep=large,row sep=large]
 K_n(D,E)\ar[r]\ar[d,hook]&
 A_n^{\Gamma_{D\times E}}\ar[d,hook,"a_n^{\Gamma_{D\times E}}"]\\
 D^{[n]}\times E^{[n]}\ar[r,"\Phi"']&
 \TV^{\Gamma_{D\times E}}.
\end{tikzcd}
\end{equation}
Thus the comparison has a constructed inverse on a canonical domain,
even before a full-domain product theorem is available.

\subsection{Maximal restrictions and composition}

The semantic spectra also allow arbitrary morphisms
$f:\overline D\to\overline E$ in $\E$, whether or not they arise
from old effective Scott maps. Write $j_{D,m}=\jsem_{D,m}$ in this
subsection.

\begin{definition}[Graded admission domain]
\label{def:graded-admission}
For $m,n\geq1$, define $\Adm_{m,n}(f)$ by
\begin{equation}\label{diag:graded-admission}
\begin{tikzcd}[column sep=large]
 \Adm_{m,n}(f)\ar[r,"\widehat f_{m,n}"]\ar[d,hook,"d_f"']&
 E^{[n]}\ar[d,hook,"j_{E,n}"]\\
 D^{[m]}\ar[r,"f j_{D,m}"']&\overline E.
\end{tikzcd}
\end{equation}
It represents a partial map $D^{[m]}\pmap E^{[n]}$ with its displayed
subobject domain. No dominance classification of that domain is
included in this definition.
\end{definition}

\begin{theorem}[Maximal restriction and substitution]
\label{thm:admissible-substitution}
The restriction in \eqref{diag:graded-admission} is maximal among
restrictions of $f j_{D,m}$ landing in $E^{[n]}$. For
$f:\overline D\to\overline E$ and
$g:\overline E\to\overline F$, the sequential domain is
\begin{equation}\label{eq:sequential-domain}
 S=\Adm_{m,n}(f)\times_{E^{[n]}}\Adm_{n,r}(g).
\end{equation}
There is a canonical mono
$S\incl\Adm_{m,r}(gf)$ on which the two restricted maps agree.
These comparisons satisfy identity and associativity coherence.
\end{theorem}
\begin{proof}
A subobject $U\incl D^{[m]}$ together with a factorization of
$f j_{D,m}|_U$ through $E^{[n]}$ is precisely a cone over the
pullback defining $\Adm_{m,n}(f)$. It therefore factors uniquely
through its domain, proving maximality.

A point of the pullback \eqref{eq:sequential-domain} is an admitted
input $x$ such that $f(x)$ is at level $n$ and $g(f(x))$ is at level
$r$. It therefore gives a point of the direct admission domain of
$gf$. Categorically the two pullback squares paste to give the unique
map $S\to\Adm_{m,r}(gf)$ over $D^{[m]}$. It is monic, since both
are subobjects there. Both output maps become $gf j_{D,m}|_S$ after
the mono $j_{F,r}$, so agree.

For three maps, either parenthesization gives the subobject on which
the first, second, and final intermediate values have their prescribed
grades. Associativity of pullbacks gives its canonical identification,
and uniqueness of the maps over the input proves the comparison
coherence. The same-grade restriction of an identity has its full
domain; its unit equations follow from the identity pullback.
\end{proof}

This is a calculus in the category of partial maps with arbitrary mono
domains. The direct admission domain may be larger than the sequential
one, because direct admission does not impose the intermediate grade.
The comparison is thus an inclusion of domains, not an unproved equality.

\begin{corollary}[Typed composition of graded points]
\label{cor:typed-composition}
There is an $\E$-enriched category with objects the pairs $(D,n)$ and
hom-objects $(E^{[r]})^{D^{[n]}}$. Its composition is evaluation and
transposition. The full-domain restrictions of the preceding theorem
compose whenever the intermediate grade matches.
\end{corollary}
\begin{proof}
Use the internal composition and identity maps of the cartesian closed
category $\E$ on the displayed objects. Their associativity and unit
equations are the exponential equations. If both restrictions have full
domain, their sequential domain is all of the input object. The preceding
theorem then makes the composite restriction full-domain as well.
\end{proof}

This typed category does not supply a shift endofunctor on all its
hom-objects. Such a shift would require an operation transporting an
arbitrary map $D^{[m]}\to E^{[n]}$ to the shifted grades and preserving
composition. The construction of that additional operation is recorded
in \cref{prob:jump-comparison}.

\section{Decisions, termination, and totality}
\label{sec:examples}

\subsection{Finite discrete targets}

\begin{theorem}[The changed Boolean target]
\label{thm:boolean-spectrum}
For the fully observed discrete object $2$ there are canonical
isomorphisms
\[
 \overline2\cong\TV,\qquad 2^{[n]}\cong\Delta_n.
\]
At level $1$ these recover $2^{[1]}\cong2$.
\end{theorem}
\begin{proof}
Here $\Gamma_2=\Sigma^2$. Let $\epsilon_0,\epsilon_1$ be the two
singleton observations. A character has values $p,q$ on them with
$p\wedge q=\bot$ and $p\vee^{\TV}q=\top$, so $q=\negT p$.
Every trace $(a,b)$ decomposes as
$\underline a\wedge\epsilon_0\vee\underline b\wedge\epsilon_1$.
Thus the character is necessarily
\begin{equation}\label{eq:boolean-character}
 \chi_p(a,b)=(p\wedge\alpha(a))\vee^{\TV}
                 (\negT p\wedge\alpha(b)).
\end{equation}
Conversely this formula is a character: it interprets a trace on a
Boolean partition, and finite distributivity and the indexed join
equations prove preservation of its operations. Evaluation at
$\epsilon_0$ is the inverse, proving $\overline2\cong\TV$.

A level-$n$ character has both $p$ and $\negT p$ in $A_n$.
Conversely these two memberships, the inclusion $A_1\incl A_n$,
and finite closure make \eqref{eq:boolean-character} jointly
$A_n$-valued. The condition is exactly $p\in A_n\cap B_n$.
This gives the second isomorphism, and \cref{prop:ambiguous-base}
gives its base instance.
\end{proof}

\begin{proposition}[Finite partition targets]
\label{prop:finite-target-products}
For a nonempty finite discrete object $F$ with all positive observations,
the object $F^{[n]}$ is the object of $A_n$-valued finite partitions
$(p_a)_{a\in F}$ of $\top$ in $\TV$. For nonempty finite discrete
$F,G$, the comparison
$(F\times G)^{[n]}\to F^{[n]}\times G^{[n]}$ is an isomorphism.
\end{proposition}
\begin{proof}
Evaluate a character on the singleton observations. Their finite
partition equations give pairwise disjoint $p_a$ with join $\top$.
Every positive trace $u$ has the finite expansion
$\bigvee_{a\in F}\underline{u(a)}\wedge\epsilon_a$, so its value is
$\bigvee_{a\in F}^{\TV}p_a\wedge\alpha(u(a))$. Conversely this
formula defines a character and is $A_n$-valued precisely when its
singleton values are. Given two partitions, the products $p_a\wedge q_b$
form a partition indexed by $F\times G$ at the same level. Taking
marginals is its inverse, by the partition equations. These operations
are morphisms by finite closure.
\end{proof}

\subsection{The intrinsic first separation}

\begin{theorem}[First separation]
\label{thm:first-separation}
For any fixed numeral $k$, the complement in $\TV$ of
$\alpha H(-,k)$ is a $B_1$-classified predicate on $P$ which is not
$A_1$-classified. In particular $\SSn_1(P)\incl\SSn_2(P)$ is proper.
\end{theorem}
\begin{proof}
Membership in $B_1$ is the definition of that marking. If the complement
also had an $A_1$-classifier, it and $H(-,k)$ would be positive
predicates with meet $\bot$ and join $\top$: their finite operations
can be checked in $\TV$ and reflected by $\alpha$. The resulting
complementary pair would give a map $P\to2$ distinguishing the
nowhere-defined function from a total constant function. This
contradicts intrinsic Rice in Paper~I. Since $B_1\incl A_2$, the
displayed predicate witnesses the proper inclusion.
\end{proof}

\begin{corollary}[A new Boolean observation]
\label{cor:new-boolean-map}
There is a map $h_k:P\to2^{[2]}$ with semantic Boolean value
$\alpha H(-,k)$. It does not factor through the natural inclusion
$2^{[1]}\incl2^{[2]}$.
\end{corollary}
\begin{proof}
Both $\alpha H(-,k)$ and its complement belong to $A_2$, by the
inclusions $A_1\incl A_2$ and $B_1\incl A_2$. Therefore
\cref{thm:boolean-spectrum} gives $h_k$. Explicitly its profile is
\eqref{eq:boolean-character} with $p=\alpha H(f,k)$, uniformly in
$f$. A factorization through $2^{[1]}\cong2$ would classify
termination by an ordinary Boolean decision, contrary to the preceding
theorem.
\end{proof}

\begin{proposition}[Totality]
\label{prop:totality-classification}
The predicate
$\Tot(f)=\forall n:\N\,iH(f,n)$ is $B_2$-classified and is not
$A_1$-classified.
\end{proposition}
\begin{proof}
Transpose $H$ to $P\to\Sigma^{\N}$ and apply
\cref{prop:universal-classifier}. Nonpositivity follows from the
intrinsic finite-information theorem of Paper~I. More explicitly, the
increasing finite partial identities have total supremum and no total
member. A positive classifier on $P$ would preserve this increasing
supremum, so could not take value $\top$ at the supremum and $\bot$
at every stage.
\end{proof}

These examples use the intrinsic computation object and its observation
theorems. No program index or universal numerical predicate has entered
their construction.

\section{Numerical normal forms and strictness}
\label{sec:numerical}

We now compare the constructed hierarchy with the numerical presentation
already available in Paper~I. This section uses its program cover
$q:\N\onto P$ and its adequate evaluation map. The previous sections
did not use this cover. A statement about external supports below is a
statement about the underlying subsets of the standard natural numbers;
it is distinguished from an assertion about all generalized predicates.

\subsection{Global presentations on partitioned domains}

\begin{proposition}[Numerical normal forms]
\label{prop:numerical-normalforms}
For each fixed $n\geq1$ and each partitioned assembly $C$, a map
$p:C\to\TV$ is $A_{n+1}$-classified if and only if there is an
actual map $v:C\times\N\to B_n$ such that
\begin{equation}\label{eq:numerical-normalform}
 p(x)=\joinT_k b_n(v(x,k)).
\end{equation}
The dual statement characterizes $B_{n+1}$ by countable meets of
$A_n$-predicates. This includes every finite power $C=\N^d$.
\end{proposition}
\begin{proof}
If $p$ has its $A_{n+1}$-classifier, apply the modest-target lifting
lemma of Paper~I to that map and the cover
$q_{n+1}:B_n^{\N}\onto A_{n+1}$. Its domain is modest by
\cref{lem:hierarchy-base} and exponential closure. The resulting
map $C\to B_n^{\N}$ transposes to the required $v$.
Conversely transpose $v$ and compose with $q_{n+1}$ to obtain the
classifier. Complement proves the dual statement.
\end{proof}

This proposition supplies an actual normal form for each numerical
predicate. It does not choose normal forms naturally for all domains,
and it does not identify the two sides of
\eqref{eq:presentation-comparison} for arbitrary $X$.

\begin{theorem}[Finite arithmetical calibration]
\label{thm:arithmetical-calibration}
For each fixed $n\geq1$ and finite $d\geq1$, the external supports
of maps $\N^d\to A_n$ are exactly the $\Sigma^0_n$ predicates, and
the external supports of maps $\N^d\to B_n$ are exactly the
$\Pi^0_n$ predicates. Here the numerical hierarchy starts with the
c.e. predicates, takes complement for $\Pi^0_n$, and forms
$\Sigma^0_{n+1}$ by existential numerical projection of uniformly
$\Pi^0_n$ predicates. The ambiguous level has support class
$\Delta^0_n$.
\end{theorem}
\begin{proof}
At level $1$, finite-hit classification on the partitioned object
$\N^d$ gives a tracked family of binary streams for any positive
predicate. Searching for a hit enumerates its support. Conversely a
c.e. predicate has a uniformly decidable bounded-stage hit test; its
binary stream and the finite-hit classifier give a map into $\Sigma$.
This is also the numerical positive-predicate comparison of Paper~I.

The object $\TV$ is the codiscrete Boolean assembly. Consequently a
map from an assembly into $\TV$ is determined by its underlying
Boolean-valued function. Semantic complement takes the complement of
the support, and semantic numerical join takes its existential
projection. Suppose the calibration holds at level $n$, for all finite
arities. By \cref{prop:numerical-normalforms}, a positive predicate
at level $n+1$ has the support of an existential projection of a
$B_n$-predicate on $\N^{d+1}$, hence a $\Sigma^0_{n+1}$ support.
Conversely a uniformly $\Pi^0_n$ predicate in those variables has a
$B_n$-classifier by the induction statement. Transpose it and apply
$q_{n+1}$ to obtain the required classifier of its projection.
Complement proves the negative assertion. Induction completes the
calibration. Finally the pullback defining $\Delta_n$ asks for both
classifiers of the same support, giving precisely the intersection
of these two support classes.
\end{proof}

The definition of the numerical hierarchy in the theorem fixes the
uniformity convention. In particular, an arbitrary external sequence
of predicates is not a numerical family unless its dependence on the
sequence variable is represented jointly. This is why the exponential
and lifting arguments are part of the comparison proof.

\subsection{Universality obtained after reconstruction}

Fix the numerical pairing isomorphism
$\langle-,-\rangle:\N\times\N\cong\N$ of Paper~I.
The first universal predicate is the existing termination observation
pulled back along the program cover:
\begin{equation}\label{eq:first-universal}
 U_1(e,x)=H(q(e),x):\N\times\N\longrightarrow A_1.
\end{equation}

\begin{lemma}[Universality at the positive base]
\label{lem:base-universal}
For every morphism $p:\N\to A_1$, there is a numeral $e$ such that
$p(x)=U_1(e,x)$ for all $x$, as maps to $A_1$.
\end{lemma}
\begin{proof}
The positive predicate defines the partial constant-zero function with
domain $p$. Its finite graph trace is positive: a finite graph belongs
precisely when each of its finitely many output values is zero and
each input satisfies $p$. The trace defines an effective point of the
partial-function domain, hence a point of $P$. Its termination
predicate is $p$. Lift that point through $q:\N\onto P$ using the
partitioned source $1$ and the modest domain $\N$. The resulting
numeral is the required $e$. Adequacy of evaluation through the
program cover identifies its termination observation with
\eqref{eq:first-universal}.
\end{proof}

\begin{theorem}[Universal numerical predicates]
\label{thm:universal-predicates}
For every fixed $n\geq1$ there is a map
$U_n:\N\times\N\to A_n$ such that every morphism $p:\N\to A_n$
is equal to $U_n(e,-)$ for some numeral $e$. The semantic values can
be chosen recursively by
\begin{equation}\label{eq:universal-successor}
 a_{n+1}U_{n+1}(e,x)=
 \joinT_k\negT\bigl(a_nU_n(e,\langle x,k\rangle)\bigr).
\end{equation}
\end{theorem}
\begin{proof}
The first level is \cref{lem:base-universal}. Given $U_n$, complement
its interpretation to obtain a joint $B_n$-predicate in $(e,x,k)$.
Transpose in $k$ and apply $q_{n+1}$. This constructs $U_{n+1}$
as an actual morphism with the stated interpretation.

Let $p:\N\to A_{n+1}$. By
\cref{prop:numerical-normalforms}, choose a map
$v:\N\times\N\to B_n$ whose join presents $p$. Complement $v$
and transport along the inverse pairing isomorphism to obtain a map
$w:\N\to A_n$. Universality at level $n$ supplies an $e$ with
$w(z)=U_n(e,z)$. Substitute $z=\langle x,k\rangle$ and use
\eqref{eq:universal-successor}; the interpreted maps $p$ and
$U_{n+1}(e,-)$ agree. Monicity of $a_{n+1}$ gives equality of the
classifiers. This proves universality at the next level.
\end{proof}

Universality here quantifies over actual morphisms $\N\to A_n$.
It does not assert, without a further parameterized enumeration proof,
that the transpose $\N\to A_n^{\N}$ is regular epic. The formulation
proved is exactly what the following diagonal argument requires.

\begin{theorem}[Strictness at every finite level]
\label{thm:finite-strictness}
For every fixed $n\geq1$, the predicate
\begin{equation}\label{eq:descriptive-diagonal}
 d_n(x)=\negT\bigl(a_nU_n(x,x)\bigr)
\end{equation}
is $B_n$-classified but is not $A_n$-classified. Consequently all
three sequences of predicate objects on $\N$ are strictly increasing:
\[
 \SSn_n(\N)\subsetneq\SSn_{n+1}(\N),\qquad
 \PPn_n(\N)\subsetneq\PPn_{n+1}(\N),\qquad
 \DDn_n(\N)\subsetneq\DDn_{n+1}(\N).
\]
\end{theorem}
\begin{proof}
The negative classifier is obtained by complement. If there were an
$A_n$-classifier for $d_n$, \cref{thm:universal-predicates} would
supply a numeral $e$ with
$a_nU_n(e,x)=d_n(x)$. At $x=e$ this gives a Boolean algebra element
equal to its own complement. Such an element would equal $\bot$ by
taking its meet with itself and then equal $\top$ by complement,
contradicting nondegeneracy of $\TV$.

Since $B_n\incl A_{n+1}$, this predicate witnesses properness of
the positive inclusion. Complement transports that properness to
the negative inclusion. Finally $d_n$ belongs to both $B_{n+1}$ and
$A_{n+1}$, using $B_n\incl B_{n+1}$ and $B_n\incl A_{n+1}$.
It therefore has a $\Delta_{n+1}$-classifier, but no
$\Delta_n$-classifier because it has no $A_n$-classifier.
The claimed proper inclusions follow.
\end{proof}

This proof derives strictness from the completed observation hierarchy
and the inherited universal computation object. It does not use an
external strictness theorem as a premise.

\section{Comparison maps and their scope}
\label{sec:comparison}

\subsection{Observation profiles and represented-space jumps}

Pauly--de~Brecht's framework represents modified open predicates by a
modified Sierpi\'nski object, and uses an admissibility comparison to
recover maps from observation profiles \cite[Sections~2.1--2.3]{paulydebrecht}.
Our corresponding object is the explicit pullback
\eqref{diag:graded-spectrum}; \cref{thm:profile-representation} gives
its reconstruction theorem. Their concrete limit-based examples
depend on representations \cite[Section~3.1]{paulydebrecht}.
De~Brecht develops jump operators as abstractions of limit computation
and relates them to levels of discontinuity \cite{debrechtjump}.

The numerical comparison is already exact by
\cref{thm:arithmetical-calibration}. A comparison on general spaces
requires more data: a correspondence between their observation
algebras, an interpretation of names, and a proof that the relevant
predicate descent and character reconstruction agree. In particular,
matching the external supports of predicates does not determine a
map between their classifier objects.

\begin{problem}[Endofunctor comparison]
\label{prob:jump-comparison}
Construct a category containing the graded effective-point objects,
with its observation structures, and construct a shift on its
hom-objects compatible with evaluation and composition. Identify its
action on the positive classifier with the successor marking, and
construct the comparison with a limit-jump presentation. The comparison
must include the resulting product maps and the admissibility domain
\eqref{eq:product-admission}.
\end{problem}

The already proved functors have type $\ObsCat\to\E$; their output
objects have not been identified with objects of $\ObsCat$. The typed
category in \cref{cor:typed-composition} supplies composition at
matching grades, and \cref{thm:admissible-substitution} supplies
maximal restrictions of semantic maps. These results give concrete
inputs to the problem, but neither supplies its shift operation.

\subsection{Boolean observation algebras and localic comparison}

Chen constructs localic descriptive notions using Boolean completions
of frames and regular completion under images \cite{chen}. This
makes comparison with our completion and descent constructions natural.
There is a distinction in the universal properties: $\Cw$ completes a
specified marking inside a fixed semantic truth object, whereas a free
Boolean completion has a universal mapping property into varying
Boolean algebras. The distinction can be stated by an explicit
evaluation map.

\begin{proposition}[Faithful semantic evaluation]
\label{prop:semantic-evaluation}
The transpose of the semantic profile defines a mono
\begin{equation}\label{eq:semantic-evaluation}
 \eta_D:\Gamma_D\incl\TV^{\overline D},\qquad
 \eta_D(u)(\chi)=\chi(u).
\end{equation}
It preserves the positive operations, interpreting them pointwise in
$\TV$.
\end{proposition}
\begin{proof}
Operation preservation is the character condition. For monicity, each
basis point $r(b)$ of Paper~I is effective and hence, by
\cref{thm:base-recovery}, supplies a point of $\overline D$ jointly
in $b$. Evaluation of $u$ there is $\alpha(u(b))$, since its positive
character evaluates the rounded trace at $b$. If two traces have the
same image under $\eta_D$, these evaluations agree for every $b$.
Monicity of $\alpha$ gives equality of the traces. This argument
retains arbitrary parameters, so proves that $\eta_D$ is monic.
\end{proof}

\begin{definition}[The semantic Boolean observation hull]
\label{def:semantic-hull}
Inside the Boolean algebra $\TV^{\overline D}$, let $\mathcal H_D$
be the intersection of all subobjects which contain the image of
$\eta_D$ and are closed under complement, finite meets, and countable
semantic joins. The intersection is taken over the power object of
$\TV^{\overline D}$.
\end{definition}

\begin{proposition}[Semantic hull and substitution]
\label{prop:semantic-hull}
The object $\mathcal H_D$ exists and has the displayed operations.
It is the least such subalgebra inside $\TV^{\overline D}$.
For an old effective map $f:D\to E$, precomposition with
$\overline f$ restricts to a map $\mathcal H_E\to\mathcal H_D$,
functorially.
\end{proposition}
\begin{proof}
The family being intersected is inhabited by the whole exponential.
Membership in every member defines a subobject using the power object.
If the inputs to any of the specified operations belong to every
member, its output belongs to every member by their closure. Thus
the operations restrict uniquely to the intersection. The universal
property of intersection proves leastness.

Precomposition preserves the pointwise operations. The inverse image
of $\mathcal H_D$ under precomposition with $\overline f$ is a
subalgebra of $\TV^{\overline E}$. It contains $\eta_E(\Gamma_E)$,
because
\[
 \eta_E(u)\,\overline f=\eta_D(f^*u).
\]
Leastness therefore makes it contain $\mathcal H_E$, giving the
restriction. Identity and composition follow from precomposition,
or from their equations after the mono into the exponential.
\end{proof}

This hull is a concrete semantic algebra to which a localic comparison
can map. Its universal property is relative to its ambient function
algebra; no freeness into arbitrary Boolean algebras is inferred.
Also, its predicates live on $\overline D$, whereas the objects
$\SSn_n(X)$ classify effective predicates on the specified input
$X$. Keeping those domains visible is necessary for a comparison
between the two constructions.

\subsection{Truth operations and local operators}

Kihara relates oracle access, changes of representation, and operations
on truth values \cite{kihara}. This is relevant to the oracle
constructions of Paper~I, but there are three different operations in
the present development. The local operator $j=\neg\neg$ supplies
the stable semantic truth object. The reflection $\Cw$ completes
markings under countable join. The operation
$M\mapsto\Cw(\negT M)$ advances the descriptive hierarchy.

Their types and universal properties distinguish them. The first acts
on $\Omega$, the second on a poset of marked subobjects of $\TV$,
and the third combines the second with an involution. By
\cref{thm:finite-strictness}, the last operation is not idempotent on
the constructed chain. It therefore cannot be identified, on that
chain, with an idempotent reflection merely from the existence of a
truth-value interpretation.

\begin{proof}[Proof of \cref{main:hierarchy,main:spectra,main:separation}]
\Cref{prop:positive-stability,thm:countable-completion-reflection,
prop:hierarchy-indexing,thm:hierarchy-operations,prop:quantifier-calculus}
give \cref{main:hierarchy}. The constructions and representation
theorem in \cref{sec:spectra}, followed by
\cref{thm:graded-functoriality,thm:semantic-products,
thm:graded-product-domain,thm:admissible-substitution}, give
\cref{main:spectra}. Finally
\cref{thm:first-separation,cor:new-boolean-map,
prop:totality-classification,thm:arithmetical-calibration,
thm:universal-predicates,thm:finite-strictness}
give \cref{main:separation}.
\end{proof}

\section{Further constructions}
\label{sec:further}

The finite theory supplies classifiers, numerical quantifiers, semantic
spectra, and exact admission domains. Several subsequent constructions
can now be formulated without choosing a syntax of descriptive codes
as their foundation.

\begin{problem}[Transfinite observation stages]
\label{prob:transfinite}
Construct an object of stage presentations and their comparison maps,
with successor given by the already constructed operation on markings.
For a limit stage, construct the family of lower-stage interpretations,
its image, and its descent theorem. Determine its numerical and spectral
operations, including the uniformity required when the lower stage
varies with an input parameter.
\end{problem}

The internal family of \cref{prop:hierarchy-indexing} provides the
finite starting diagram. A transfinite construction must additionally
exhibit the stage comparisons and recursion principle; a family of
externally specified ordinals alone would not provide those arrows.

\begin{problem}[Analytic projection]
\label{prob:analytic}
Starting from the inherited object of infinite numerical data, construct
the image of the semantic existential projection on the appropriate
relation object. Identify its presentation covers and the descent
needed for a classifier on arbitrary input objects. Construct the
corresponding spectral interpretation and its substitution maps.
\end{problem}

The operation $j\exists$ on all stable predicates is already available,
but its restriction to an effective marking is a separate image
construction. \Cref{thm:predicate-descent} applies once that
interpretation and its image have been constructed; it does not
identify the resulting effective class in advance.

\begin{problem}[Well-founded relations and ranks]
\label{prob:ranks}
Construct the relevant objects of well-founded relations, their
recursion operations, and rank comparisons. Establish the relation
between inductive accessibility and the negative assertion that no
infinite descending sequence exists. Derive the boundedness statements
needed to compare their stages with the descriptive hierarchy.
\end{problem}

These are distinct constructive notions of well-foundedness until the
comparison has been proved. The stable semantic interpretation alone
does not identify them. The rank and boundedness constructions are also
the point at which analogues of admissibility and stronger ordinal
closure properties can be investigated within the framework.

\begin{problem}[Structural set-theoretic interpretation]
\label{prob:structural-sets}
From the resulting well-founded structures, construct a category of
families and a distinguished class of maps, together with the generic
family and the closure, collection, and descent constructions required
for an algebraic set-theoretic interpretation. Construct the comparison
between material descriptions and structural families and prove its
coherence with the descriptive operations.
\end{problem}

The regular descent and spectral reconstruction proved here specify
how descriptions can acquire invariant denotations. They do not by
themselves construct the small-map structure or a material universe.
Those structures are the outputs requested by
\cref{prob:structural-sets}, after the analytic and rank constructions
have been supplied.

\appendix
\section{Interface with the preceding manuscript}
\label{app:imports}

This manuscript is self-contained as a LaTeX source. Mathematical
dependencies on Paper~I are cited by result name, so that recompiling
either paper does not invalidate the references. The following table
records the exact labels in the supplied source of Paper~I for future
cross-document linking. These strings are identifiers, not additional
mathematical premises.

\begingroup\small
\begin{longtable}{L{.32\linewidth}L{.40\linewidth}L{.19\linewidth}}
\toprule
Result in Paper~I & Source label & Use here\\
\midrule
\endfirsthead
\toprule
Result in Paper~I & Source label & Use here\\
\midrule
\endhead
\bottomrule
\endfoot
Indexed coreflection of admissible families &
\texttt{thm:admissibility-coreflection} & Maximal domains\\
Modest-target lifting & \texttt{lem:modest-target-lift} & Cover lifting\\
Function names and subobjects & \texttt{lem:modestclosure} & Modest closure\\
Finite-hit classification & \texttt{thm:finitehit} & Positive base\\
Relative foundation & \texttt{thm:relativefoundation} & Numerical joins\\
Finite-hit name presentations & \texttt{lem:sigmamodest} & Modesty of $\Sigma$\\
Two-spectrum reconstruction & \texttt{jd:thm:spectra} & Base recovery\\
Positive observation duality & \texttt{jd:cor:duality} & Maps on profiles\\
Products of searchable presentations & \texttt{prop:productpresentation} & Product traces\\
Intrinsic Rice and halting & \texttt{thm:intrinsichalting} & First separation\\
Intrinsic finite information & \texttt{thm:observationfullness} & Totality\\
Nonpositivity corollary & \texttt{cor:intrinsicnonpositive} & Separation\\
Numerical program cover & \texttt{prop:codecover} & Universality\\
Modal-singleton classifier for sheafification &
\texttt{thm:oracle-singleton-classifier} & Semantic background\\
\end{longtable}
\endgroup

\section{Notation and logical interpretations}
\label{app:notation}

\begin{longtable}{L{.26\linewidth}L{.66\linewidth}}
\toprule
Notation & Meaning\\
\midrule
\endfirsthead
\toprule
Notation & Meaning\\
\midrule
\endhead
\bottomrule
\endfoot
$\E$ & The effective topos of Paper~I.\\
$i:\Sigma\incl\Omega$ & The inherited positive classifier.\\
$\TV=\Omega_{\neg\neg}$ & Stable semantic truth; its joins are
double-negation closures of ambient joins.\\
$\alpha:\Sigma\incl\TV$ & The stable interpretation of positive truth.\\
$\Cw$ & Countable-join reflection on modest markings of $\TV$.\\
$A_n,B_n$ & Positive and negative markings at finite level $n$.\\
$\Delta_n$ & Their intersection over $\TV$, classifying ambiguous predicates.\\
$\SSn_n(X),\PPn_n(X)$ & The exponential predicate objects $A_n^X,B_n^X$.\\
$\Gamma_D$ & The positive trace algebra inherited from the searchable domain $D$.\\
$\overline D$ & The semantic spectrum $\Spec_{\TV}(\Gamma_D)$.\\
$D^{[n]}$ & The subobject of semantic characters with jointly $A_n$-valued profile.\\
$\Adm_{m,n}(f)$ & The maximal level-$m$ domain on which a semantic map
has level-$n$ output.\\
$\mathcal H_D$ & The least Boolean, countably join-closed semantic
observation subalgebra containing $\Gamma_D$.\\
\end{longtable}

In particular, $\Sigma$ is not the subobject classifier of a new ambient
topos, $A_n$ is not asserted to be a dominance, and $\Delta_n$ does not
classify ordinary complementary subobjects unless the complement and
join conditions descend to $\Omega$. At level $1$ that descent is
proved in \cref{prop:ambiguous-base}. Likewise the letters
$\Sigma^0_n,\Pi^0_n,\Delta^0_n$ are used only in the numerical support
comparison of \cref{thm:arithmetical-calibration}, where their meaning
and equivalence have been established.

\begin{thebibliography}{99}

\bibitem{sheridan}
Conor Sheridan.
\newblock \emph{Effective Domains, Oracles, and Forcing in the Effective Topos}.
\newblock Manuscript, November 2026. Cited as Paper~I.

\bibitem{hyland}
J. M. E. Hyland.
\newblock The effective topos.
\newblock In A. S. Troelstra and D. van Dalen, editors,
\emph{The L. E. J. Brouwer Centenary Symposium}, Studies in Logic and the
Foundations of Mathematics 110, pages 165--216. North-Holland, 1982.

\bibitem{vanoosten}
Jaap van Oosten.
\newblock \emph{Realizability: An Introduction to its Categorical Side}.
\newblock Studies in Logic and the Foundations of Mathematics 152.
Elsevier, 2008.

\bibitem{hylandvariations}
J. M. E. Hyland.
\newblock Variations on realizability: realizing the propositional axiom of choice.
\newblock Author manuscript dated August 28, 2003.
\newblock \url{https://www.dpmms.cam.ac.uk/~jmeh1/Research/Publications/2002/vor02.pdf}.

\bibitem{paulydebrecht}
Arno Pauly and Matthew de Brecht.
\newblock Towards synthetic descriptive set theory: an instantiation with represented spaces.
\newblock arXiv:1307.1850, version 2, 2014.
\newblock \url{https://arxiv.org/abs/1307.1850}.

\bibitem{debrechtjump}
Matthew de Brecht.
\newblock Levels of discontinuity, limit-computability, and jump operators.
\newblock arXiv:1312.0697, 2013.
\newblock \url{https://arxiv.org/abs/1312.0697}.

\bibitem{chen}
Ruiyuan Chen.
\newblock Borel and analytic sets in locales.
\newblock arXiv:2011.00437, 2020.
\newblock \url{https://arxiv.org/abs/2011.00437}.

\bibitem{kihara}
Takayuki Kihara.
\newblock Rethinking the notion of oracle: a prequel to Lawvere--Tierney topologies for computability theorists.
\newblock arXiv:2202.00188, version 4, 2024.
\newblock \url{https://arxiv.org/abs/2202.00188}.

\end{thebibliography}
\end{document}


